\documentclass[10pt,a4paper]{amsart}
\usepackage[margin=19mm]{geometry}
\usepackage{amsmath,amssymb,mathtools}
\usepackage[scr=rsfs,cal=euler]{mathalfa}
\usepackage{xfrac}
\usepackage[unicode,hidelinks,bookmarks=false]{hyperref}
\newcommand{\ie}{i.e.\ }
\newcommand{\cf}{cf.\ }
\newcommand{\resp}{respectively\ }
\newcommand{\Subsection}{\subsection}
\providecommand{\nobreakdash}{\nobreak\hbox{-}}
\setlength{\emergencystretch}{2em}
\multlinegap0pt
\usepackage{bm}
\usepackage{bbm}
\usepackage{dsfont}
\usepackage{xspace}
\usepackage{cancel}
\usepackage{mathtools}
\usepackage{amsthm}
\makeatletter
\@ifundefined{theorem}{\newtheorem{theorem}{Theorem}}{}
\@ifundefined{lemma}{\newtheorem{lemma}[theorem]{Lemma}}{}
\@ifundefined{proposition}{\newtheorem{proposition}[theorem]{Proposition}}{}
\@ifundefined{remark}{\newtheorem{remark}[theorem]{Remark}}{}
\makeatother
\DeclareMathOperator{\supp}{supp}
\newcommand{\N}{\mathbb{N}}
\newcommand{\R}{\mathbb{R}}
\newcommand{\Tree}{\mathcal{T}}
\newcommand{\s}{\sigma}
\title[Embedding $\ell\_2$ and $J$ into subspaces of $JT$ and $JT^*$]{Embedding $\ell\_2$ and $J$ into subspaces of $JT$ and $JT^*$}
\author{Spiros A. Argyros, Manuel Gonzalez,, Pavlos Motakis}
\date{Source: arXiv:2603.17886v2 (CC BY 4.0)}
\begin{document}
\maketitle

\noindent\emph{Source and licence.} Embedding $\ell_2$ and $J$ into subspaces of $JT$ and $JT^*$. Version arXiv:2603.17886v2 (CC BY 4.0), distributed under the Creative Commons Attribution 4.0 International licence, as recorded in the arXiv licence metadata for this exact version and shown on the abstract page https://arxiv.org/abs/2603.17886v2. Full rights notice: licenses/arxiv-2603.17886-CC-BY-4.0.txt.\par\smallskip

\noindent\textit{Editorial scope.} Counted unit: the Proposition whose source label is \texttt{P11} in the article (source0.tex lines 1004--1008, source label \texttt{P11}), delivered together with the article's own complete original proof of it (source0.tex lines 1009--1087, one \texttt{proof} environment of 79 lines). No mathematics is written, repaired or re-derived: statement and proof are the article's own text line for line. Required context retained uncounted: P\_A (lines 238--240); A-I (lines 247--249); P10 (lines 832--840); compl.J (lines 923--924); L8 (lines 989--993); R1 (lines 995--1002). Every definition, display and label of those lines is retained unchanged. Omitted from this package and not redistributed: 1--237, 241--246, 250--831, 841--922, 925--988, 994--994, 1003--1003, 1088--1201. Nothing inside a retained excerpt is omitted or altered: every retained body is delivered exactly as the article's lines are written, so no omission receipt of any kind accompanies this package. Citations: the retained excerpts carry no citation command at all, so no bibliography entry of the article is transcribed and no reference is redistributed. Reference adaptation: none was needed and none was made; every reference inside the delivered unit resolves inside it, so no unresolved reference or citation is produced. Printed slips kept literal: no printed word, symbol, hypothesis or number of the counted statement or of its proof was altered. Layout and class: the article's own class is \texttt{amsart} with packages including amscd, amssymb, verbatim, color; the wrapper is the lane's pinned \texttt{amsart} editorial wrapper at 10pt on A4, which restates the theorem-like environments the retained text needs and adds the article's own macro definitions that the retained text needs (\texttt{\textbackslash N}, \texttt{\textbackslash R}, \texttt{\textbackslash Tree}, \texttt{\textbackslash s}, \texttt{\textbackslash supp}), with extra wrapper packages (bm, bbm, dsfont, xspace, cancel, mathtools). The delivered numbering is the unit's own. Assets: no retained span contains a figure environment, a graphics-inclusion command, a PDF-page inclusion, a table or a TikZ picture; no image file is copied into this package, the package declares no asset dependency and no image file is redistributed with it. Licence: version arXiv:2603.17886v2 of this article is distributed under the Creative Commons Attribution 4.0 International licence, as recorded in the arXiv licence metadata for this exact version and shown on the abstract page https://arxiv.org/abs/2603.17886v2. A full rights notice is delivered at licenses/arxiv-2603.17886-CC-BY-4.0.txt. Characters: every retained excerpt is pure ASCII, so the delivered engine, XeLaTeX with its default Latin Modern text font, reports no missing character.
\begin{proposition}\label{P_A}
Let $A\subset \Tree$ such that for each segment $s$ in $\Tree$ the intersection $s\cap A$ is again a segment. Then $P_A(\sum_{k=1}^\infty a_k e_{k}) = \sum_{k=1; k\in A}^\infty a_k e_{k}$ defines a norm-one projection $P_A:JT\to JT$.
\end{proposition}

\begin{theorem}\label{A-I}
If $(x_n)_n$ is a seminormalized block sequence in $JT$ such that $\sigma^*(x_n) \to 0$ for every $\sigma \in \Gamma$, then there exists a subsequence $(x_{n_k})_k$ of $(x_n)_n$ equivalent to the unit vector basis of $\ell_2$. In particular, $(x_n)_n$ is a weakly null sequence.
\end{theorem}

\begin{proposition}\label{P10}
Let $(x_n)_n$ be a block sequence in $B_J$ such that
\[ \lim_n I_\infty^*(x_n) = \alpha \neq 0 \quad \text{and} \quad \lim_n \|x_n\| =\beta \geq 0,\]
and let also $0< \epsilon < 1$. Then there exists a subsequence $(x_{n_k})_k$ such that for every $(\lambda_k)_k$ in $\R$ we have
\begin{equation}\label{eqn:eq.6}
(1-\epsilon) |\alpha|\Big\| \sum_{k=1}^\infty \lambda_k e_k \Big\| \leq \Big\| \sum_{k=1}^\infty \lambda_k x_{n_k} \Big\| \leq 3(1+\epsilon)(|\alpha|+ \beta)\Big\| \sum_{k=1}^\infty \lambda_k e_k \Big\|.
\end{equation}
Hence $(x_{n_k})_k$ is equivalent to the basis $(e_k)_k$ of $J$.
\end{proposition}

\begin{lemma}\label{compl.J} Let $(x_n)_n$ be a block sequence in $J$ equivalent to the basis of $J$. Then the space generated by some subsequence of $(x_n)_n$ is complemented in $J$.
\end{lemma}

\begin{lemma}\label{L8}
Suppose that $A = \{ \s_{i,>m} : i\in F\}$ is incomparable. Then for each $x\in JT$,
\[ \|P_A(x)\| = \Big( \sum_{\s \in A} \|P_{\s_{> m}}(x)\|_J^2 \Big)^{1/2}, \]
where $\|\cdot\|_J$ denotes the norm of the James space $J$.
\end{lemma}

\begin{remark}\label{R1}
Let $x^{**} = \sum_{i=1}^\infty a_i\s_i^{**}$ with  $\|x^{**}\| = 1$, and let $(x_n)_n$ be a block sequence such that $x_n \xrightarrow{w^*} x^{**}$. In the next proposition, we will assume that $(x_n)_n$ satisfies the following additional properties:
\begin{enumerate}
    \item[i)] $\lim_n \|x_n\| = b \geq 1$.
    \item[ii)] For every $i \in \N$, $\lim_n \|P_{\s_i}(x_n)\| = b_i \geq |a_i|$.
\end{enumerate}
Clearly, every $(x_n)_n$ with $x_n \xrightarrow{w^*} x^{**}$ has a subsequence satisfying (i) and (ii). Moreover, it is easy to check that $(\sum_{i=1}^\infty b_i^2)^{1/2} \leq b$.
\end{remark}

\begin{proposition}\label{P11}
Let $(x_n)_n$ be a block sequence in $JT$ such that
$x_n \xrightarrow{w^*} x^{**}$ with $x^{**} = \sum_{i=1}^\infty a_i \s_i^{**}$ and $\|x^{**}\| = 1$.
Then $(x_n)_n$ has a subsequence equivalent to the basis of $J$ generating a complemented subspace of $JT$.
\end{proposition}

\begin{proof}
As we have noticed, we can assume that $(x_n)_n$ satisfies (i) and (ii) in Remark \ref{R1}. For $0< \epsilon < 1$, we choose a sequence $(\delta_k)_k$ of positive reals such that
\begin{equation}\label{eqn:eq.13}
\delta_1 = b \quad \text{and} \quad \sum_{k=1}^\infty \delta_k < (1 +\epsilon) b.
\end{equation}

Next we partition $\N$ into successive disjoint intervals $(F_k)_k$ satisfying for $k\geq 2$:
\begin{equation}\label{eqn:eq.12}
3(1+\epsilon) \Big(\sum_{i\in F_k}(|a_i| + b_i)^2 \Big)^{1/2} < \delta_k,
\end{equation}
and we choose an increasing sequence of naturals $(m_k)_k$ such that the set
\[ A_k =\Big\{ \s_{i,> m_k} : i \in F_k\Big\} \]
consists of incomparable final segments. Proposition \ref{P_A} yields that $\|P_{A_k}\| = 1$.
By induction, we define a decreasing sequence $(L_k)_k$ of infinite subsets of $\N$ such that the following facts are satisfied. (In the following $(e_n)_n$ denotes the basis of $J$.)

\begin{equation}\label{eqn:eq.11}
\begin{split}
(i) \quad &\forall n\in L_k, \text{ we have } m_k < \supp(x_n) \text{ and } \forall i\in F_k, \quad \Big| \s_i^*(x_n) -a_i \Big| < \epsilon |a_i| \text{ and }  \\
&\quad \Big| \|P_{\s_i}(x_n)\| - b_i \Big| < \epsilon b_i. \\
(ii) \quad &\forall i\in F_k \text{ and scalars }(\lambda_n)_{n\in F_k},\\
&\quad (1-\epsilon) |a_i| \Big\| \sum_{n\in L_k}\lambda_ne_n \Big\|_J\leq \Big\| \sum_{n\in L_k} \lambda_nP_{\s_i}(x_n) \Big\|  \leq  3(1+\epsilon)(|a_i|+ b_i)\Big\| \sum_{n \in L_k} \lambda_ne_n \Big\|. \\
(iii) \quad &\text{For } k\geq 2, \quad \Big\| \sum_{n\in L_k} \lambda_nP_{A_k}(x_n) \Big\| \leq 3(1+\epsilon)\Big(\sum_{i\in F_k}(|a_i|+ b_i)^2\Big)^{1/2}\Big\| \sum_{n \in L_k} \lambda_ne_n \Big\| _J\leq \\ &\quad \leq  \delta_k \Big\| \sum_{n \in L_k} \lambda_ne_n \Big\|_J.
\end{split}
\end{equation}

The inductive definition of $L_k$ satisfying (i) follows from the fact that $(x_n)_n$ is a block sequence, $\|P_{\s_i}(x_n)\| \to b_i$ and $\s_i^*(x_n) \to a_i$.

Property (ii) follows from Proposition \ref{P10} and Property (iii) follows from (ii), Lemma \ref{L8} and the properties of $(\delta_k)_k$ (\ref{eqn:eq.12}).

This completes the inductive definition of the sequence $(L_k)_k$.

Choose an increasing sequence $(n_k)_k$ with $n_k\in L_k$ and set
\[ v_k = P_{\cup_{p\leq k} A_p}(x_{n_k}) \quad \text{and} \quad w_k = x_{n_k} - v_k. \]
We will show the following facts:

1) The block sequence $(v_k)_k$ is dominated by the basis of $J$.

2) The sequence $(w_k)_k$ is either norm null or has a subsequence equivalent to the $\ell_2$ basis.

Notice that if (1) and (2) hold, then there exists a subsequence of $(x_{n_k})_k$ which is equivalent to the basis of $J$. If $(w_k)_k$ is norm null, then this follows from a standard perturbation argument. 
Otherwise, $(w_k)$ is equivalent to the $\ell_2$ basis. Then there exist $c_1, c_2$ positive constants such that
\begin{equation}\label{eqn:eq.10}
\begin{split}
\Big\|\sum_{k=1}^\infty\lambda_k x_{n_k} \Big\| &\leq \Big\| \sum_{k=1}^\infty\lambda_k v_k \Big\| + \Big\| \sum_{k=1}^\infty\lambda_k w_k \Big\| \\
&\leq c_1 \Big\| \sum_{k=1}^\infty\lambda_k e_k \Big\|_J + c_2\Big( \sum_{k=1}^\infty\lambda_k^2 \Big)^{1/2} \leq (c_1+c_2)\Big\| \sum_{k=1}^\infty\lambda_k e_k \Big\|_J.
\end{split}
\end{equation}
The last inequality follows from the fact  that the $\ell_2$ norm is dominated by the $J$ norm in the basis $(e_k)_k$ of $J$. If a subsequence of $(w_k)_k$ is equivalent to the $\ell_2$ basis, then (\ref{eqn:eq.10}) remains valid for this subsequence. Also, if $(w_k)_k$ is norm null, then (1) and a permutation argument yield a subsequence of $(x_{n_k})_k$ dominated by the basis of $J$.

For the opposite inequality, namely the basis of $J$ is dominated by $(x_{n_k})_n$, we notice that $G = \{n_k : k\in\N\} \subset L_1$. Hence (\ref{eqn:eq.11}) (ii) gives:
\[ (1-\epsilon) |a_1| \Big\| \sum_{n\in G}\lambda_ne_n \Big\|_J \leq \Big\| \sum_{n\in G} \lambda_nP_{\s_1}(x_n) \Big\| \leq \Big\| \sum_{n\in G} \lambda_nx_n \Big\|. \]
This ends the proof of the equivalence. It remains to prove (1) and (2).

To prove (1), let $(\lambda_k)_{k\leq m} \subset \R$. By the definition of $(v_k)_k$ we have
\begin{equation}\label{eqn:eq.14}
\sum_{k=1}^m \lambda_k v_k = \sum_{k=1}^m \lambda_k P_{A_1}(v_k) + \sum_{k=2}^m \lambda_k P_{A_2}(v_k) + \dots + \lambda_m P_{A_m}(v_m).
\end{equation}
Since for every $p\geq k$ we have $P_{A_k}(v_p) = P_{A_k}(x_{n_p})$ and $(n_p)_{p\geq k} \subset L_k$ for $1 < k\leq m$, from the bimonotonicity of the basis of $J$ and (\ref{eqn:eq.11}) (iii) we derive:
\begin{equation}\label{eqn:eq.15}
\Big\| \sum_{p=k}^m \lambda_pP_{A_k}(v_p) \Big\| \leq \delta_k \Big\| \sum_{p=k}^m \lambda_p e_p \Big\|_J \leq \delta_k \Big\| \sum_{p=1}^m \lambda_p e_p \Big\|_J.
\end{equation}
For $k=1$ we have
\begin{equation}\label{eqn:eq.17}
\Big\| \sum_{p=1}^m \lambda_pP_{A_1}(v_p) \Big\| \leq 3\Big( \sum_{i\in F_1} (|a_i| + b_i)^2\Big)^{1/2} \Big\| \sum_{p=1}^m \lambda_p e_p \Big\|_J \leq 3\sqrt{2} (1+ \epsilon) b\Big\| \sum_{p=1}^m \lambda_p e_p \Big\|_J.
\end{equation}
Using the triangle inequality, (\ref{eqn:eq.13}) and (\ref{eqn:eq.17}) we get
\begin{equation}\label{eqn:eq.16}
\Big\|\sum_{k=1}^m \lambda_k v_k \Big\| < 6(1+\epsilon)b \Big\|\sum_{k=1}^m \lambda_k e_k \Big\|_J.
\end{equation}
This completes the proof of (1).

To prove (2), we show that $v_k \xrightarrow{w^*} x^{**}$. Indeed, if $\s = \s_i$, then $\s_i^*(v_k) \to a_i = x^{**}(\s_i^*)$. Assume $\s\neq \s_i$ for all $i \in \N$.
We assert that $\lim_k \s^*(v_k) = 0$. Observe from (\ref{eqn:eq.11}) (i) we have that for every $\delta > 0$ there exists a finite initial segment $F$ of $\N$ and $k_0$ such that for $k>k_0$ and $\s \neq \s_i$ for $i\in F$ we have $\| P_\s(v_k)\| < \delta$.

Therefore the assertion is true and $v_k \xrightarrow{w^*} x^{**}$.
Since $x_{n_k} \xrightarrow{w^*} x^{**}$ and $w_k = x_{n_k} - v_k$, we conclude that $w_k \xrightarrow{w} 0$. This and Theorem \ref{A-I} proves (2) and completes the proof that the subsequence $(x_{n_k})_k$ is equivalent to the basis of $J$.

For the complementation of $[(x_{n_k})_k]$ we observe that $P_{\s_1}$ (defined above) is an isomorphism on $[(x_{n_k})_k]$. So, Lemma \ref{compl.J} yields the property.
\end{proof}

\end{document}
